\subsection{Credential preparation and Secure Messaging}
\label{sec:credential-sm-flow}

OSDP Secure Channel protects communication between the ACU and PD.
Credential Secure Messaging (SM) protects credential commands and replies between
the PD and the credential. Credential SM can operate over contact or contactless.
The Virtual Contact Interface (VCI) allows a contactless credential to perform
operations that otherwise require contact, subject to the credential's access rules.
When the credential requires pairing, the PD establishes SM before submitting the
eight-digit pairing code. Pairing is the default, recommended configuration in SP~800-73-5 Part~1,
Section~5.5. The pairing code and the cardholder's PIN serve separate purposes:
pairing enables VCI; PIN verification authorizes PIN-protected operations.

\paragraph{Credential presentation}
The PD \textbf{SHALL} select the configured application on an interface allowed
by \texttt{osdp\_PIVMODE}. For a PIV application, the PD \textbf{SHALL} inspect
the SELECT response's application property template (tag \code{0x61}). Its
cryptographic algorithm template (tag \code{0xAC}) contains algorithm identifiers
in tag \code{0x80}. The PD \textbf{SHALL} examine all such entries. An identifier
of \code{0x27} or \code{0x2E} advertises credential SM; SP~800-73-5 Part~2,
Section~3.1.1, requires this indication when the PIV application supports SM.
The PD \textbf{SHALL} distinguish a malformed response from a valid response
that advertises neither SM algorithm.

\paragraph{Discovery Object}
The Discovery Object (\code{0x7E}) describes the credential's PIN and VCI policy.
Its PIN Usage Policy (tag \code{0x5F2F}) identifies the supported PIN options,
the preferred PIN when both PIN types are enabled, and whether VCI requires
pairing. The PD uses this information to prepare PIN and VCI operations.
SP~800-73-5 Part~1, Section~3.3.2, defines the policy bytes.

The PD \textbf{SHALL} read the Discovery Object when preparing a credential-side PIN
operation or attempting SM setup, unless the PD already holds the policy for
the selected application on the current credential. Automatic reads \textbf{SHALL} be limited to those
workflows. Collecting a PIN for later use requires no credential read. A status query
reports the policy already read, or \code{0xFFFF} when unavailable. The ACU can
also request the object explicitly with \texttt{osdp\_PIVGETDATA}.

In the first policy byte, mask \code{0x08} indicates VCI support. When VCI is
supported, mask \code{0x04} permits VCI without a pairing code; otherwise pairing
is required. If the object is absent, VCI is unavailable. The SELECT response
remains the source of SM capability information.

\paragraph{Automatic SM preparation}
After selecting a PIV application, the PD checks the requirements for SM setup.
If PIVMODE flag \code{0x02} is set, the PD supports an SM algorithm advertised
by the credential, and at least one trust anchor is loaded, the PD \textbf{SHALL}
attempt SM setup. This applies to both contact and contactless credentials, in
commanded mode and PIV Auto.

Selecting a different application starts this check again. Selecting the same
PIV application again preserves the current SM session and PIN-verification
status. Reselection alone therefore requires neither new SM setup nor repeated
PIN verification. The credential's access rules still apply to each operation. These selection
rules follow SP~800-73-5 Part~2, Section~3.1.1.

The PD's trust store contains the anchors loaded through
\texttt{osdp\_PIVVCILOADTA}. The PD checks its own store before starting SM.
An empty store prevents setup. If anchors are present, the PD must still find
one that can validate this credential's certificate, directly or through an
intermediate certificate. Loading an anchor makes it available for the next
attempt. To start a new attempt with a credential already present, the ACU reapplies
PIVMODE; the mode-change clearing rules apply.

The credential returns its SM card-verifiable certificate (CVC) during key
establishment. The PD \textbf{SHALL} validate that certificate using a loaded
anchor, directly or through an intermediate CVC, and verify the authentication
cryptogram as specified in
SP~800-73-5 Part~2, Section~4.1. The PD \textbf{SHALL} report SM as established
only after those checks succeed.

\paragraph{Proceeding without SM}
When SM is disabled, unsupported, or unavailable because a trust anchor is
missing, the PD \textbf{SHALL} check whether the operation requested by the
ACU or configured in PIV Auto requires SM. If SM is required, the PD
\textbf{SHALL} terminate the operation without sending the operation's command
to the credential and return the defined failure reply. Otherwise,
the PD \textbf{SHALL} perform the requested operation under the credential's access
rules. The PD \textbf{SHALL NOT} substitute another key, algorithm, interface,
or operation to avoid an SM requirement.

PIV key \code{0x9A} requires PIN verification or a supported biometric
comparison within the credential, as specified in SP~800-73-5 Part~1,
Section~5.1, Table~5. The PIN-based contactless PIV Auto workflow with key \code{0x9A}
therefore requires SM, VCI, and successful PIN verification. The PD
\textbf{SHALL} check those requirements in that order before sending the
authentication command. Missing PIN verification returns Result \code{0x05},
Detail \code{0x02} (PIN verification required).
If SM is unavailable, the PD reports \texttt{osdp\_PIVSTATUS} Result
\code{0x02} (SM failure) with the applicable detail code. PIV Auto with key
\code{0x9E} can continue without SM and reports the credential-authentication result.
A commanded PIN verification over contact can also proceed without SM.

The PD may start key establishment before discovering that the required
trust anchor is missing. In that case, the PD \textbf{SHALL} discard any
session keys derived during setup. The PD then handles the requested operation
using the rules above.

In commanded mode, the ACU obtains the preparation result with
\texttt{osdp\_CARDSTATUS}. The PD reports the established SM and VCI state in
\texttt{osdp\_CARDSTATUSR}. A missing anchor leaves permitted data reads
available. For example, the ACU can use \texttt{osdp\_PIVGETDATA} to read the
Secure Messaging Certificate Signer object (\code{0x5FC122}) and inspect the
signer certificate and any intermediate CVC. The credential's own SM CVC is returned
during key establishment; that CVC is separate from the signer data object.
The ACU can request \texttt{osdp\_CARDSTATUS} Status Type \code{0x02} to obtain
the setup outcome and the hash of the CVC received during that attempt, including when
setup failed.

\paragraph{Failed security checks}
A missing or invalid credential CVC, invalid authentication cryptogram, or
failure to verify an SM reply \textbf{SHALL} end the current operation with a
failure result. The PD
\textbf{SHALL} block further data access and authentication commands until the
credential is removed, a different application is successfully selected, or
PIVMODE is reapplied. PIVSELECT remains available to select a different
application; selecting the same application preserves the blocked state.
This prevents continued use of a session that failed its security checks.
The rule also applies when automatic SM setup fails before an ACU command
arrives. PIVMODE, trust-store management, and SM diagnostic requests remain available.
For a standard status request (Status Type \code{0x01}), the PD terminates the
request without sending further commands or data to the credential and returns
\texttt{osdp\_PIVERROR} \code{0x1023} (credential security conditions not
satisfied). This distinguishes a blocked credential from operation without SM.

For blocked operations, the PD reports \texttt{osdp\_PIVSPER} \code{0x08}
(security condition not met) for SPE, or \texttt{osdp\_PIVERROR} \code{0x1023}
(credential security conditions not satisfied) for commands that use PIVERROR.

If an operation requires SM, VCI, or PIN verification and that requirement is
unmet, the PD \textbf{SHALL} report failure in the operation's completion reply.
PIV Auto uses \texttt{osdp\_PIVSTATUS}; SPE uses \texttt{osdp\_PIVSPER}.

Figure~\ref{fig:sm-preparation} summarizes credential preparation and the
decision to proceed with the requested operation.
\begin{figure}[htbp]
\centering
\begin{tikzpicture}[x=1cm,y=1cm,font=\sffamily\small,
  box/.style={draw=OSDPBlue,rounded corners=2pt,fill=OSDPLightGray,
    align=center,text width=8.2cm,inner sep=6pt},
  side/.style={draw=OSDPBlue,align=center,text width=3.3cm,inner sep=5pt},
  flow/.style={->,draw=OSDPBlue,thick}]
\node[box] (select) at (0,0) {Select the PIV application\\Check the SELECT reply for SM support};
\node[box] (ready) at (0,-1.5) {SM enabled, supported, and a trust anchor loaded?};
\node[box] (sm) at (0,-3) {Read Discovery Object for PIN and pairing policy\\Establish SM: check the certificate and cryptogram};
\node[box] (access) at (0,-5) {Check the requested operation's access requirements\\For a PIN operation, read Discovery policy if needed};
\node[side] (done) at (-2.5,-7) {Conditions satisfied\\Perform operation};
\node[side] (fail) at (2.5,-7) {Report failure};
\draw[flow] (select) -- (ready);
\draw[flow] (ready) -- node[right]{Yes} (sm);
\draw[flow] (ready.west) -- node[above,font=\scriptsize]{No} (-5.4,-1.5)
  -- (-5.4,-5) -- (access.west);
\draw[flow] (sm) -- node[right,align=left,font=\scriptsize]{SM established, or\\required trust anchor missing} (access);
\draw[flow] (sm.east) -- (5.4,-3) -- (5.4,-7) -- (fail.east);
\node[align=center,font=\scriptsize,text width=1.8cm,fill=white,inner sep=1pt]
  at (5.4,-4) {Security check\\failed, or\\SM error};
\draw[flow] (access.south) -- (0,-5.8) -- (-2.5,-5.8) -- (done.north);
\draw[flow] (access.south) -- (0,-5.8) -- (2.5,-5.8) -- (fail.north)
  node[midway,right,font=\scriptsize]{Conditions unmet};
\end{tikzpicture}
\caption{How the PD prepares a credential. The PD establishes SM when enabled and
supported, using a loaded trust anchor. The credential's access rules determine which
operations can proceed. Contactless PIN verification requires VCI.}
\label{fig:sm-preparation}
\end{figure}


\paragraph{Contactless operation}
When pairing is required, the ACU sends \code{osdp_PIVXMITPAIRING}
after SM is established. The PD submits the pairing
code to the credential through SM. A successful verification enables VCI. For a credential
whose Discovery Object permits VCI without pairing, establishing SM also
enables VCI.
PIN verification over contactless requires VCI. Biometric comparison within the credential
and pairing-code submission can use
contactless SM without VCI, subject to the credential's access rules.
SM and VCI preserve the credential's operation-specific restrictions, including the
contact-interface restriction on ordinary PIV PUT DATA.

\paragraph{Protected credential commands}
When credential SM is established, the PD \textbf{SHALL} use that SM session
to protect PIV commands that support SM. The PD \textbf{SHALL} verify and
unwrap each protected reply before interpreting the credential data or status.
These requirements apply to data access, authentication,
PIN verification, and biometric operations. The PD \textbf{SHALL} follow
SP~800-73-5 Part~2, Section~4.2, for encryption, MACs, counters, and credential-command
and response chaining. The PD's GET RESPONSE exchanges follow that
section's chaining rules.

The PD \textbf{SHALL} discard the SM session keys on credential reset, an SM
error, or replacement by newly established session keys, as specified in
SP~800-73-5 Part~2, Section~4.3. This requirement applies to the outer SW
processing status defined in Section~4.2.7. The protected command status in
tag \code{0x99} reports the result of the credential operation separately.
When SM ends, the PD \textbf{SHALL}
clear dependent VCI and PIN-verification status and apply the PIN-cache
clearing rules.

A verified protected reply reporting a credential rejection \textbf{SHALL}
end the requested operation with that failure while preserving SM. The PD
\textbf{SHALL} update verification status according to the credential's
response. An incorrect PIN (\code{0x63Cx}) clears PIN-verification status
and the cached PIN; established pairing and VCI remain available. An incorrect
pairing code (\code{0x6300}) clears pairing verification, making VCI unavailable
when pairing is required. These credential rules are defined in SP~800-73-5
Part~2, Sections~3.2.1.1 and~3.2.1.3. The SM diagnostic outcome remains
unchanged. A failed certificate, cryptogram, or protected-reply check follows
the blocked-state rules above.

After an authenticated PIN rejection, the ACU can query remaining attempts
with PIVSPE Mode \code{0x06}, acquire a corrected PIN, and request verification
again. After a pairing rejection, the ACU can submit a corrected pairing code.
Neither retry requires a new SM session.

After an SM failure, the ACU \textbf{SHALL} reapply PIVMODE
to start a new SM attempt before retrying a protected operation.
For a commanded retry, the ACU disables PIV Auto
in that configuration, waits for SM establishment, supplies any required
pairing code, and requests PIN acquisition and verification again. The PD
\textbf{SHALL NOT} retry VERIFY automatically.
Figures~\ref{fig:sm-pin-rejection} and~\ref{fig:sm-pairing-rejection} show
the ACU and PD exchanges after authenticated rejection.
Figure~\ref{fig:sm-error-recovery} shows recovery after SM failure.
The PD \textbf{SHALL} keep a completed operation's reply until it is delivered
or canceled.

\begin{figure}[htbp]
\centering
\begin{tikzpicture}[x=1cm,y=1cm,font=\sffamily\small,
  actor/.style={draw=OSDPBlue,fill=OSDPLightGray,rounded corners=3pt,
    minimum width=2.3cm,minimum height=0.65cm},
  message/.style={->,thick,draw=OSDPBlue},
  reply/.style={message,dashed},
  note/.style={draw=OSDPBlue,fill=OSDPLightGray,align=center,
    text width=6cm,inner sep=7pt}]
\node[actor] (acu) at (0,0) {ACU};
\node[actor] (pd) at (5.7,0) {PD};
\node[actor] (credential) at (11.4,0) {Credential};
\foreach \x in {0,5.7,11.4} \draw[gray,dashed] (\x,-0.4) -- (\x,-11.0);
\draw[message] (0,-1.1) -- node[above,align=center]{PIVSPE Mode 0x02\\Verify the cached PIN} (5.7,-1.1);
\draw[reply] (5.7,-1.9) -- node[above]{ACK: request accepted} (0,-1.9);
\draw[message] (5.7,-3.0) -- node[above,align=center]{VERIFY PIN (key 0x80)\\Protected by existing SM} (11.4,-3.0);
\draw[reply] (11.4,-4.1) -- node[above,align=center]{Verified protected result: 0x63Cx\\Incorrect PIN; x attempts remain} (5.7,-4.1);
\node[note] at (5.7,-5.25) {PD erases the cached PIN and clears\\PIN-verification status.\\SM, pairing, and VCI remain established.};
\draw[message] (0,-7.2) -- node[above]{POLL} (5.7,-7.2);
\draw[reply] (5.7,-8.5) -- node[above,align=center]{PIVSPER 0x03\\PIN mismatch; request ended} (0,-8.5);
\draw[message] (0,-10.0) -- node[above,align=center]{ACU chooses to retry:\\PIVSPE Mode 0x04, acquire a new PIN} (5.7,-10.0);
\node[note,text width=11.4cm] at (5.7,-12.0) {The ACU can query remaining attempts with Mode 0x06.
After new PIN acquisition completes, a new Mode 0x02 request uses the existing SM and VCI.};
\end{tikzpicture}
\caption{Incorrect PIN over established VCI. The PD verifies the protected reply
before reporting the rejection. PIN entry occurs at the PD; the ACU never receives
the PIN. Acquisition and its completion follow the normal ACK/POLL exchange.}
\label{fig:sm-pin-rejection}
\end{figure}

\begin{figure}[htbp]
\centering
\begin{tikzpicture}[x=1cm,y=1cm,font=\sffamily\small,
  actor/.style={draw=OSDPBlue,fill=OSDPLightGray,rounded corners=3pt,
    minimum width=2.3cm,minimum height=0.65cm},
  message/.style={->,thick,draw=OSDPBlue},
  reply/.style={message,dashed},
  note/.style={draw=OSDPBlue,fill=OSDPLightGray,align=center,
    text width=6cm,inner sep=7pt}]
\node[actor] at (0,0) {ACU};
\node[actor] at (5.7,0) {PD};
\node[actor] at (11.4,0) {Credential};
\foreach \x in {0,5.7,11.4} \draw[gray,dashed] (\x,-0.4) -- (\x,-11.0);
\draw[message] (0,-1.1) -- node[above,align=center]{PIVXMITPAIRING\\Pairing code supplied by the ACU} (5.7,-1.1);
\draw[reply] (5.7,-1.9) -- node[above]{ACK: request accepted} (0,-1.9);
\draw[message] (5.7,-3.0) -- node[above,align=center]{VERIFY pairing code (key 0x98)\\Protected by existing SM} (11.4,-3.0);
\draw[reply] (11.4,-4.1) -- node[above,align=center]{Verified protected result: 0x6300\\Incorrect pairing code} (5.7,-4.1);
\node[note] at (5.7,-5.7) {Pairing verification is cleared.\\SM remains established.\\VCI still requires successful pairing.};
\draw[message] (0,-7.5) -- node[above]{POLL} (5.7,-7.5);
\draw[reply] (5.7,-9.4) -- node[above,align=center,font=\sffamily\footnotesize]{PIVERROR 0x1023: security condition unmet\\SW1/SW2 = 0x6300\\Incorrect pairing code} (0,-9.4);
\draw[message] (0,-10.8) -- node[above,align=center]{ACU chooses to retry:\\PIVXMITPAIRING with a corrected code} (5.7,-10.8);
\node[note,text width=11.4cm] at (5.7,-12.5) {The new request uses the existing SM session.
VCI becomes available after successful pairing; the PD reports that state in CARDSTATUSR.};
\end{tikzpicture}
\caption{Incorrect pairing code on a contactless credential that requires pairing.
The pairing code originates at the ACU. Each retry is a new command; the PD
never repeats VERIFY automatically.}
\label{fig:sm-pairing-rejection}
\end{figure}

\begin{figure}[htbp]
\centering
\begin{tikzpicture}[font=\sffamily,
  panel/.style={draw=OSDPBlue,rounded corners=3pt,fill=OSDPLightGray,
    text width=12cm,align=left,inner sep=9pt},
  flow/.style={->,draw=OSDPBlue,thick}]
\node[panel] (failure) {\textbf{Input: SM failure or an unconfirmed VERIFY exchange}\\[3pt]
Clear dependent VCI, PIN-verification, and cached-PIN state. Report failure once.
Block credential data and authentication commands. Diagnostics and recovery
controls remain available. An unconfirmed PIN VERIFY returns PIVSPER 0x0B
(credential communication failed or timed out). A pairing VERIFY timeout
returns PIVERROR 0x0002 (smart-card timeout). Both require preparation to restart.};
\node[panel,anchor=north] (diagnostic) at ([yshift=-0.5cm]failure.south)
{\textbf{Output to the ACU: failed operation and diagnostics}\\[3pt]
ACU sends POLL; PD returns the operation's failure reply.
ACU then sends CARDSTATUS Type 0x02; PD acknowledges and returns
CARDSTATUSR Type 0x02 fragments with the SM outcome and CVC hash on later polls.};
\node[panel,anchor=north] (restart) at ([yshift=-0.5cm]diagnostic.south)
{\textbf{New input from the ACU: restart preparation}\\[3pt]
ACU corrects the cause and reapplies PIVMODE with SM enabled and PIV Auto
disabled. PD starts SM setup. ACU requests CARDSTATUS Type 0x01 and polls
for CARDSTATUSR Type 0x01 to learn whether SM was established.};
\node[panel,anchor=north] (retry) at ([yshift=-0.5cm]restart.south)
{\textbf{Resume the requested operation}\\[3pt]
After SM succeeds, supply a pairing code when required. Acquire and verify a
PIN if the operation requires one. Each VERIFY requires a new ACU request.};
\draw[flow] (failure.south) -- (diagnostic.north);
\draw[flow] (diagnostic.south) -- (restart.north);
\draw[flow] (restart.south) -- (retry.north);
\end{tikzpicture}
\caption{Recovery after SM failure or an unconfirmed VERIFY. The ACU receives
the failed operation's result before requesting diagnostics and restarting preparation.}
\label{fig:sm-error-recovery}
\end{figure}


\paragraph{Commanded PIN verification}
For a PIN entered before credential presentation, the ACU first requests Mode
\code{0x04} (acquire and cache) and waits for completion. After credential presentation,
the PD prepares the credential. The ACU checks credential status and supplies any required
pairing code before requesting cached-PIN verification.
Figure~\ref{fig:sm-commanded-sequence} shows a successful contactless workflow
with SM enabled, a suitable trust anchor loaded, and PIV Auto disabled.
The ACU obtains the pairing code during enrollment or through another
deployment-defined process. On contact, pairing is omitted
and Mode \code{0x03} uses the selected contact interface, with SM when established.

\begin{figure}[htbp]
\centering
\begin{tikzpicture}[x=1cm,y=1cm,font=\sffamily,
  stage/.style={draw=OSDPBlue,rounded corners=3pt,fill=OSDPLightGray,
    text width=12.5cm,align=left,inner sep=10pt},
  flow/.style={->,draw=OSDPBlue,thick}]
\node[stage] (cache) at (0,0) {
  \textbf{1. Collect the cardholder's PIN}\\[4pt]
  The ACU asks the PD to collect and cache the PIN. The cardholder enters
  the PIN at the reader. The PD retains the PIN and reports completion
  to the ACU.\\[4pt]
  {\small Command: PIVSPE Mode 0x04. Completion: PIVSPER.}};
\node[stage,anchor=north] (prepare) at ([yshift=-0.55cm]cache.south) {
  \textbf{2. Prepare the contactless credential}\\[4pt]
  When the credential is presented, the PD selects the PIV application and checks
  SM support. The PD reads the Discovery Object for PIN and pairing policy,
  then establishes SM using a loaded trust anchor. The ACU requests credential
  status and learns that SM is established and pairing is required.\\[4pt]
  {\small Command: CARDSTATUS Type 0x01. Completion: CARDSTATUSR Type 0x01.}};
\node[stage,anchor=north] (pair) at ([yshift=-0.55cm]prepare.south) {
  \textbf{3. Enable VCI with the pairing code}\\[4pt]
  The ACU supplies the credential's eight-digit pairing code. The PD sends the code
  to the credential through SM. The credential verifies the code, enabling VCI for
  contactless PIN verification. The PD reports the updated status.\\[4pt]
  {\small Command: PIVXMITPAIRING. Completion: CARDSTATUSR Type 0x01.}};
\node[stage,anchor=north] (verify) at ([yshift=-0.55cm]pair.south) {
  \textbf{4. Verify the cached PIN}\\[4pt]
  The ACU requests PIN verification. The PD sends the cached PIN to the credential
  through SM. The credential checks the PIN, and the PD reports the result to the
  ACU. A successful result permits operations protected by that PIN, subject
  to the credential's access rules.\\[4pt]
  {\small Command: PIVSPE Mode 0x02. Completion: PIVSPER.}};
\draw[flow] (cache.south) -- (prepare.north);
\draw[flow] (prepare.south) -- (pair.north);
\draw[flow] (pair.south) -- (verify.north);
\end{tikzpicture}
\caption{Contactless PIN verification with a pairing code. The ACU directs the
reader (PD); the credential verifies the pairing code and PIN. Each OSDP command
receives ACK when accepted and a completion reply on a later POLL.
Appendix~\ref{app:sm-vci-examples} shows the packet exchanges.}
\label{fig:sm-commanded-sequence}
\end{figure}


\paragraph{ACU recovery after a PIV Auto failure}
PIV Auto reports one result for the configured authentication attempt. The ACU
can then request SM diagnostics, correct missing trust material, or choose
another authentication workflow. Loading an anchor makes that anchor available
to the PD. Reapplying PIVMODE starts a new automatic attempt and clears the
previous PIN and SM state. The ACU \textbf{SHALL} finish trust recovery before
pairing and PIN verification that depend on the new SM session.

When SM is established but the configured authentication requires unsupported
VCI, the PD \textbf{SHALL} terminate the attempt and report
\texttt{osdp\_PIVSTATUS} Result \code{0x05}, Detail \code{0x03} (required VCI
unsupported). The PD sends no authentication command for that attempt.
Missing SM support is reported as Result \code{0x02}, Detail \code{0x05}
(required SM unsupported). A trust-anchor update cannot add either capability.
The ACU must select a supported workflow or reject the transaction.
Appendix~\ref{ex:sm-recovery} shows trust recovery followed by commanded
pairing, PIN verification, and authentication.
